% ch3.tex -- Module 5, Chapter 3: Circular Motion: Angles, Energy, Work and Torque
% Dot points: M5.IQ2.04 (total energy and work in UCM), M5.IQ2.05 (torque),
% plus angular velocity from M5.IQ2.03 (omega = Delta theta / t).
\chapter{Circular Motion: Angles, Energy, Work and Torque}\label{ch3:chap}
\covers{M5.IQ2.04, M5.IQ2.05}

\inquiry{Why do objects move in circles?}

\begin{goals}
  \item measure angles in radians and use $\omega = \dfrac{\Delta\theta}{t}$ and $v = \omega r$
  \item explain why the centripetal force does no work, so the speed stays constant
  \item explain how work done by other forces changes the total energy
  \item use energy to find speeds and forces in vertical circles
  \item calculate torque with $\tau = r_\perp F = rF\sin\theta$ and explain how torque changes rotation
\end{goals}

\begin{syllabus}
  \item M5.IQ2.03: solve problems with $\omega = \Delta\theta/t$ (angular velocity; the other
        circular-motion formulas are in \cref{ch2:chap})
  \item M5.IQ2.04: investigate how the total energy of an object in uniform circular motion is
        related to the work done on it
  \item M5.IQ2.05: investigate how the rotation of mechanical systems depends on the torque
        applied, using $\tau = r_\perp F = rF\sin\theta$
\end{syllabus}

%% ==================================================================
\section{Angles and angular velocity}\label{ch3:sec:omega}
\covers{M5.IQ2.03}

\begin{bigidea}
Every point on a spinning object turns through the same angle in the same time. So every
point has the same angular velocity $\omega$. Points further from the centre travel further
in each turn. So they move faster: $v = \omega r$.
\end{bigidea}

\begin{keywords}
  \kw{Radian (rad)}{a unit for angles. One full turn is $2\pi$ rad, the same as \ang{360}.}
  \kw{Angular displacement ($\Delta\theta$)}{the angle an object turns through, measured at the centre.}
  \kw{Angular velocity ($\omega$)}{the angle turned each second. Unit: \si{rad.s^{-1}}.}
\end{keywords}

Think of a merry-go-round. Mia sits near the centre and Sam sits at the edge. They go round
once in the same time. But Sam travels around a bigger circle, so he moves faster. The
merry-go-round turns at the same rate for both of them. We describe this turning with
angles. The \term{angular displacement} ($\Delta\theta$) is the angle that a line from the
centre to the object turns through. ($\theta$, ``theta'', is the Greek letter for an angle.)

\begin{toolkit}{Radians and turns}
\begin{sidebyside}[0.68]
For $\omega$ and for arc length, angles must be in \term{radians} (rad). We use radians
because $s = r\theta$ and $v = \omega r$ are only true when the angle is in radians. One
radian is the angle at the centre when the arc is as long as the radius (\cref{ch3:fig:radian}).
\begin{itemize}
  \item One full turn $= 2\pi \text{ rad} = \ang{360}$, so $1 \text{ rad} \approx \ang{57.3}$.
        Degrees to radians: multiply by $\dfrac{\pi}{180}$ (so $\ang{90} = \qty{1.57}{rad}$).
  \item Turns to radians: multiply by $2\pi$ (3 turns $\approx \qty{18.8}{rad}$).
        \term{rpm} (turns per minute) to \si{rad.s^{-1}}: multiply by $\dfrac{2\pi}{60}$.
        \si{rad.s^{-1}} to rpm: multiply by $\dfrac{60}{2\pi}$.
  \item \term{Arc length} (distance along the edge): $s = r\theta$, with $\theta$ in rad (not on
        the sheet). On the formula sheet $s$ means displacement; here $s$ means the distance
        along the arc.
  \item $\num{7.27e-5}$ means $7.27 \div 100\,000 = 0.0000727$. Type it with the EXP
        (or $\times10^x$) key.
\end{itemize}
\nextside
  \centering
  \begin{tikzpicture}[x=1cm, y=1cm]
    \def\R{1.35}
    \draw[hscmuted!60, line width=0.8pt] (0,0) circle (\R);
    \draw[line width=2.4pt, hscvel] (0:\R) arc (0:57.2958:\R);
    \draw[hscink, line width=0.9pt] (0,0) -- node[below, diagramlabel] {$r$} (0:\R);
    \draw[hscink, line width=0.9pt] (0,0) -- node[above left, diagramlabel, inner sep=1pt] {$r$} (57.2958:\R);
    \draw[hscink, line width=0.6pt] (0:0.4) arc (0:57.2958:0.4);
    \node[diagramlabel, anchor=west, font=\footnotesize] at (24:0.42) {1 rad};
    \node[diagramlabel, text=hscvel, anchor=west, align=left] at (30:\R+0.08) {arc\\$= r$};
    \fill[hscink] (0,0) circle (1.5pt);
  \end{tikzpicture}
  \captionof{figure}{One radian.}\label{ch3:fig:radian}
\end{sidebyside}
\end{toolkit}

\begin{sidebyside}[0.47]
  \centering
  \begin{tikzpicture}[x=1cm, y=1cm]
    \def\R{1.8}
    \fill[hscbrand!6] (0,0) circle (\R);
    \draw[hscmuted!70, line width=0.8pt] (0,0) circle (\R);
    \draw[guide] (0,0) -- (-30:\R);
    \draw[guide] (0,0) -- (40:\R);
    \draw[displacement] (-30:0.85) arc (-30:40:0.85);
    \draw[displacement] (-30:1.7) arc (-30:40:1.7);
    \draw[hscink, line width=0.6pt] (-30:0.35) arc (-30:40:0.35);
    \node[diagramlabel, font=\footnotesize] at (5:0.6) {$\Delta\theta$};
    \node[diagramlabel, text=hscmuted] at (5:1.25) {$s_\mathrm{A}$};
    \node[diagramlabel, text=hscmuted, anchor=west] at (5:1.8) {$s_\mathrm{B}$};
    \coordinate (A) at (40:0.85);
    \coordinate (B) at (40:1.7);
    \draw[velocity] (A) -- ++(130:0.55) node[left, text=hscvel, inner sep=1pt] {$v_\mathrm{A}$};
    \draw[velocity] (B) -- ++(130:1.1) node[left, text=hscvel, inner sep=1pt] {$v_\mathrm{B}$};
    \fill[object] (A) circle (3.5pt);
    \fill[object] (B) circle (3.5pt);
    \node[diagramlabel, anchor=south, inner sep=1pt] at ($(A)+(80:0.2)$) {A};
    \node[diagramlabel, anchor=south, inner sep=1pt] at ($(B)+(80:0.2)$) {B};
    \fill[hscink] (0,0) circle (1.5pt);
  \end{tikzpicture}
  \captionof{figure}{A and B turn through the same angle $\Delta\theta$ in the same time.}\label{ch3:fig:sameomega}
\nextside
  The \term{angular velocity} ($\omega$, ``omega'') is the angle turned each second.
  \begin{itemize}
    \item In \cref{ch3:fig:sameomega}, A and B turn through the same angle in the same time.
          They have the same $\omega$.
    \item B is twice as far from the centre. Its arc $s_\mathrm{B}$ is twice as long as
          $s_\mathrm{A}$, so it moves twice as fast.
    \item The speed of a point along its circle is its \term{tangential speed} $v$.
  \end{itemize}
\end{sidebyside}

\begin{formulacard}[unbreakable]{Angular velocity}
  \cardformula{\omega = \frac{\Delta\theta}{t}}
  \begin{symbols}
    \sym{$\omega$}{angular velocity: angle turned each second}{\si{rad.s^{-1}}}
    \sym{$\Delta\theta$}{angular displacement: angle turned}{\si{rad}}
    \sym{$t$}{time taken}{\si{s}}
  \end{symbols}
  \inwords{angular velocity = angle turned $\div$ time taken.}
  \whentouse{anything that spins or moves round a circle: wheels, fans, the Earth.}
  \whennottouse{the angle is in degrees or turns. Change it to radians first.}
  \rearranged{$\Delta\theta = \omega t$ \qquad $t = \dfrac{\Delta\theta}{\omega}$}
  \tensecond{\qty{12}{rad} in \qty{3.0}{s}: $\omega = \dfrac{12}{3.0} = \qty{4.0}{rad.s^{-1}}$.}
  \trap{one turn is $2\pi$ rad, not 1 rad. ``5 turns'' means $\Delta\theta = 10\pi$ rad.}
\end{formulacard}

\begin{formulacard}[notonsheet]{Angular velocity from the period}
  \cardformula{\omega = \frac{2\pi}{T} = 2\pi f}
  \begin{symbols}
    \sym{$\omega$}{angular velocity}{\si{rad.s^{-1}}}
    \sym{$T$}{period: time for one full turn}{\si{s}}
    \sym{$f$}{frequency: turns each second}{\si{Hz} ($= \si{s^{-1}}$)}
  \end{symbols}
  \inwords{one full turn ($2\pi$ rad) takes one period $T$.}
  \whentouse{you know the period $T$ or the frequency $f$. For rpm, first find
    $f = \text{rpm} \div 60$ (turns each second).}
  \whennottouse{the turning rate changes. Then it gives only an average $\omega$.}
  \rearranged{$T = \dfrac{2\pi}{\omega}$ \qquad $f = \dfrac{\omega}{2\pi}$}
  \tensecond{the Earth turns once a day ($T = \qty{24}{h} = \qty{86400}{s}$, the value HSC
    questions use), so $\omega = \dfrac{2\pi}{86400} = \qty{7.27e-5}{rad.s^{-1}}$.}
  \derive{start from formulas on the sheet.
    \begin{align*}
      \omega &= \frac{\Delta\theta}{t} = \frac{2\pi}{T}
        && \text{one turn: $\Delta\theta = 2\pi$ in the time $t = T$}\\
      \omega &= 2\pi \times \frac{1}{T} = 2\pi f
        && \text{$f = \dfrac{1}{T}$ (on the sheet)}
    \end{align*}}
  \trap{$T$ must be in seconds: 1 min = \qty{60}{s}, 1 h = \qty{3600}{s}.}
\end{formulacard}

\begin{formulacard}[notonsheet]{Speed from angular velocity}
  \cardformula{v = \omega r}
  \begin{symbols}
    \sym{$v$}{tangential speed of a point}{\si{m.s^{-1}}}
    \sym{$\omega$}{angular velocity}{\si{rad.s^{-1}}}
    \sym{$r$}{distance of the point from the centre (axis)}{\si{m}}
  \end{symbols}
  \inwords{for the same $\omega$, a point twice as far out moves twice as fast.}
  \whentouse{you know $\omega$ and need the speed of one point on a spinning object (or the reverse).}
  \whennottouse{$\omega$ is in rpm or degrees per second. Change it to \si{rad.s^{-1}} first.}
  \rearranged{$\omega = \dfrac{v}{r}$ \qquad $r = \dfrac{v}{\omega}$}
  \tensecond{at the equator, $v = \omega r_\mathrm{E}$, where $r_\mathrm{E} = \qty{6.371e6}{m}$ is
    the radius of the Earth (data sheet): $v = \num{7.27e-5} \times \num{6.371e6} = \qty{463}{m.s^{-1}}$.}
  \derive{start from $v = \dfrac{2\pi r}{T}$ (on the sheet).
    \begin{align*}
      v &= \frac{2\pi r}{T} = \left(\frac{2\pi}{T}\right) r = \omega r
        && \text{group $\dfrac{2\pi}{T}$, which is $\omega$}\\
      a_\mathrm{c} &= \frac{v^2}{r} = \frac{(\omega r)^2}{r} = \omega^2 r
        && \text{put $v = \omega r$ into $a_\mathrm{c} = \dfrac{v^2}{r}$; cancel one $r$}
    \end{align*}}
  \trap{measure $r$ from the centre: the radius, not the diameter.}
\end{formulacard}

\begin{workedexample}{A ceiling fan}\label{ch3:we:fan}
A ceiling fan makes 15 full turns in \qty{6.0}{s} at a steady rate. The tip of each blade
is \qty{0.60}{m} from the centre. Calculate (a) the angular velocity and (b) the speed of a
blade tip.
\given{15 turns in $t = \qty{6.0}{s}$, \ $r = \qty{0.60}{m}$}
\find{(a) $\omega$ \quad (b) the tip speed $v$}
\formula{$\omega = \dfrac{\Delta\theta}{t}$ \quad then \quad $v = \omega r$}
\substitute{\eqlines{\Delta\theta &= 15 \times 2\pi = \qty{94.248}{rad} \qquad \text{(turns to radians)}\\
  \omega &= \qty{94.248}{rad} \div \qty{6.0}{s} = \qty{15.708}{rad.s^{-1}} \qquad \text{(angle $\div$ time)}\\
  v &= \qty{15.708}{rad.s^{-1}} \times \qty{0.60}{m} = \qty{9.425}{m.s^{-1}} \qquad \text{(speed $= \omega \times r$)}}}
\answer{(a) \final{$\omega = \qty{16}{rad.s^{-1}}$} \quad (b) \final{$v = \qty{9.4}{m.s^{-1}}$}
  (2 significant figures, like the data)}
\checkstep{another route: $T = 6.0 \div 15 = \qty{0.40}{s}$ and
  $v = 2\pi r \div T = 2\pi \times 0.60 \div 0.40 = \qty{9.42}{m.s^{-1}}$ \cmark. That is about
  \qty{34}{km.h^{-1}}, a sensible speed for a fan tip.}
\end{workedexample}

\begin{workedexample}{A laboratory centrifuge}\label{ch3:we:washer}
A laboratory centrifuge spins at 3600 rpm. A test tube sits \qty{0.0800}{m} from the axis.
Find (a) the angular velocity, (b) the speed of the tube and (c) its centripetal
acceleration. Compare (c) with $g$.
\given{3600 turns in \qty{60}{s}, \ $r = \qty{0.0800}{m}$, \ $g = \qty{9.8}{m.s^{-2}}$}
\find{(a) $\omega$ \quad (b) $v$ \quad (c) $a_\mathrm{c}$}
\formula{$\omega = \dfrac{\Delta\theta}{t}$, \quad $v = \omega r$, \quad
  $a_\mathrm{c} = \dfrac{v^2}{r}$ (on the sheet, \cref{ch2:chap})}
\substitute{\eqlines{\omega &= (3600 \times 2\pi \text{ rad}) \div \qty{60}{s}
    = \qty{376.99}{rad.s^{-1}} \qquad \text{(rpm to \si{rad.s^{-1}})}\\
  v &= \qty{376.99}{rad.s^{-1}} \times \qty{0.0800}{m} = \qty{30.159}{m.s^{-1}} \qquad \text{($v = \omega r$)}\\
  a_\mathrm{c} &= (\qty{30.159}{m.s^{-1}})^2 \div \qty{0.0800}{m} = \qty{11370}{m.s^{-2}} \qquad
    \text{($a_\mathrm{c} = v^2 \div r$)}}}
\answer{(a) \final{$\omega = \qty{377}{rad.s^{-1}}$} \quad (b) \final{$v = \qty{30.2}{m.s^{-1}}$}
  \quad (c) \final{$a_\mathrm{c} = \qty{1.14e4}{m.s^{-2}}$}, about $11\,370 \div 9.8 \approx 1200$ times $g$.
  (Keep the unrounded values in your calculator until you round here.)}
\checkstep{$\omega^2 r = 376.99^2 \times 0.0800 = \qty{1.14e4}{m.s^{-2}}$, the same \cmark. The tube
  must be pushed towards the centre with a force about 1200 times its weight. So centrifuge
  tubes must be strong.}
\end{workedexample}

\begin{mistakes}
\begin{itemize}
  \item \xmark\ Putting turns, degrees or rpm straight into $\omega = \dfrac{\Delta\theta}{t}$.
        \cmark\ Change to radians: 1 turn $= 2\pi$ rad.
  \item \xmark\ ``All points on a wheel move at the same speed.''
        \cmark\ They have the same $\omega$. The speed $v = \omega r$ grows with $r$.
  \item \xmark\ Using the diameter in $v = \omega r$. \cmark\ $r$ is measured from the centre.
\end{itemize}
\end{mistakes}

\begin{examtip}
Multiple-choice questions often ask what is the same for every point on a spinning object.
The answer: $\omega$, $T$ and $f$ are the same, but $v$ and $a_\mathrm{c}$ grow with $r$. In calculations, write the
conversion line, for example $3600 \text{ rpm} = 3600 \times 2\pi \div 60 =
\qty{377}{rad.s^{-1}}$. It often earns a mark.
\end{examtip}

\begin{quickcheck}
  \item Change to radians: (a) \ang{60} \quad (b) 2.5 turns. (c) Change 3000 rpm to \si{rad.s^{-1}}.
  \item A turntable spins at 45 rpm. Point A is \qty{0.050}{m} and point B is \qty{0.15}{m} from
        the centre. (a) Which point has the larger $\omega$? (b) Find the speed of each point.
  \item A car travels at \qty{72}{km.h^{-1}}. Its wheels have a radius of \qty{0.30}{m} and do not
        slip. Find the angular velocity of a wheel in \si{rad.s^{-1}} and in rpm.
\end{quickcheck}

%% ==================================================================
\section{Work and energy in circular motion}\label{ch3:sec:work}
\covers{M5.IQ2.04}

\begin{bigidea}
A force does work only when part of it points along the motion. In uniform circular motion
the net force is at \ang{90} to the motion. So it does no work, and the speed stays constant.
If only gravity does work in a vertical circle (no friction, no motor), the speed changes,
but $K + U$ stays constant.
\end{bigidea}

\begin{keywords}
  \kw{Work ($W$)}{energy transferred by a force when the object moves along the direction of
    the force (or of part of it). Unit: joule (J).}
  \kw{Total mechanical energy ($E$)}{kinetic energy plus potential energy: $E = K + U$.}
  \kw{Non-uniform circular motion}{moving in a circle with a changing speed.}
  \kw{Apparent weight}{the push of the seat or floor on you. It is what you feel as your weight.}
  \kw{Dip / crest}{the low point of a valley / the top of a hill in a track.}
\end{keywords}

\begin{refresher}{Work and energy (Year 11)}
\begin{itemize}
  \item \term{Kinetic energy} (energy of motion) is $K = \tfrac{1}{2}mv^2$.
        \term{Gravitational potential energy} ($U$) is energy stored because of height. When the
        height changes by $\Delta h$, it changes by $\Delta U = mg\Delta h$.
  \item A force $F$ at angle $\theta$ to the motion has a part $F\cos\theta$ along the motion.
        Only this part does work.
  \item \term{Work--energy theorem}: the work done by the net force equals the change in
        kinetic energy, $W_\mathrm{net} = \Delta K$ (not on the sheet).
\end{itemize}
\end{refresher}

Imagine you carry a heavy bag across a flat room. Your hand pushes \emph{up}, but the bag moves
\emph{sideways}. Your push is at \ang{90} to the motion, so it does no work on the bag.

The same thing happens in a circle (\cref{ch3:fig:work}a). The \term{centripetal force} (the net
force towards the centre) is at \ang{90} to the velocity at every instant.

\begin{figure}
  \centering
  \begin{tikzpicture}[x=1cm, y=1cm]
    % ---- (a) uniform
    \begin{scope}
      \def\R{1.2}
      \draw[hscmuted!60, line width=0.8pt] (0,0) circle (\R);
      \fill[hscink] (0,0) circle (1.5pt);
      \node[diagramlabel, left, inner sep=3pt, font=\footnotesize] at (0,0) {centre};
      \coordinate (P) at (\R,0);
      \draw[hscink, line width=0.5pt] ($(P)+(-0.22,0)$) -- ($(P)+(-0.22,0.22)$) -- ($(P)+(0,0.22)$);
      \draw[force] (P) -- ++(-0.8,0) node[pos=0.5, below, inner sep=3pt, text=hscforce] {$F_\mathrm{c}$};
      \draw[velocity] (P) -- ++(0,1.0) node[right, inner sep=2pt, text=hscvel] {$v$};
      \draw[displacement] ($(P)+(0.3,0.0)$) -- ++(0,0.45) node[right, inner sep=2pt, text=hscmuted, font=\footnotesize] {$\Delta s$};
      \fill[object] (P) circle (4.5pt);
      \node[diagramlabel, align=center] at (0.55,-1.6) {\textbf{(a) Uniform:} $W = F_\mathrm{c}\,\Delta s\cos\ang{90} = 0$};
    \end{scope}
    % ---- (b) speeding up
    \begin{scope}[xshift=6.9cm]
      \def\R{1.2}
      \draw[hscmuted!60, line width=0.8pt] (0,0) circle (\R);
      \fill[hscink] (0,0) circle (1.5pt);
      \node[diagramlabel, left, inner sep=3pt, font=\footnotesize] at (0,0) {centre};
      \coordinate (P) at (\R,0);
      \draw[guide] ($(P)+(-0.8,0)$) -- ($(P)+(-0.8,0.6)$) -- ($(P)+(0,0.6)$);
      \draw[force, dashed, line width=1pt] (P) -- ++(-0.8,0) node[pos=0.78, below, inner sep=4pt, text=hscforce, font=\footnotesize] {radial};
      \draw[force, dashed, line width=1pt] (P) -- ++(0,0.6) node[right, inner sep=3pt, text=hscforce, font=\footnotesize, align=left] {tangential:\\does work};
      \draw[force] (P) -- ++(-0.8,0.6) node[above left, inner sep=1pt, text=hscforce] {$F_\mathrm{net}$};
      \draw[->, hscmuted, line width=0.6pt] (-55:1.5) arc (-55:-15:1.5);
      \node[diagramlabel, text=hscmuted, anchor=west, font=\footnotesize] at (-35:1.55) {motion};
      \fill[object] (P) circle (4.5pt);
      \node[diagramlabel, align=center] at (0.55,-1.6) {\textbf{(b) Speeding up}};
    \end{scope}
  \end{tikzpicture}
  \caption{(a) The net force is at \ang{90} to the velocity: no work.
    (b) The net force has a part along the motion.}\label{ch3:fig:work}
\end{figure}

Split the path into tiny steps $\Delta s$ along the tangent. For each step,
$W = F_\mathrm{c}\,\Delta s\cos\ang{90} = 0$. So the net work is zero.

By the work--energy theorem, $\Delta K = W_\mathrm{net} = 0$. The kinetic energy and the speed stay
constant. The centripetal force changes only the \emph{direction} of the velocity.

If the object speeds up, the net force has two parts: a \emph{radial} part (towards the
centre) and a \emph{tangential} part (along the path). Only the tangential part does work
(\cref{ch3:fig:work}b).

\begin{formulacard}{Work done by a force}
  \cardformula{W = F_\parallel s = Fs\cos\theta}
  \begin{symbols}
    \sym{$W$}{work done by the force}{\si{J}}
    \sym{$F$, $F_\parallel$}{the force, and its part along the displacement}{\si{N}}
    \sym{$s$}{displacement of the object (for a force that always points along a curved path:
      the distance along the path)}{\si{m}}
    \sym{$\theta$}{angle between the force and the displacement}{\si{\degree} (calculator in
      degree mode)}
  \end{symbols}
  \inwords{only the part of the force along the motion does work.}
  \whentouse{a steady force acts on an object while it moves. For a circle, use it for each
    tiny step along the tangent.}
  \whennottouse{to find the work of a centripetal force by multiplying it by the distance
    round the circle. It is at \ang{90} to every step, so $W = 0$.}
  \rearranged{$F = \dfrac{W}{s\cos\theta}$ \qquad $s = \dfrac{W}{F\cos\theta}$}
  \tensecond{\qty{20}{N} at \ang{60} to a \qty{5.0}{m} displacement:
    $W = 20 \times 5.0 \times \cos\ang{60} = \qty{50}{J}$. At \ang{90}, $W = 0$.}
  \trap{a ball whirled on a \qty{4.3}{N} string through half a turn (\qty{2.4}{m}) does
    \emph{not} gain $4.3 \times 2.4 = \qty{10}{J}$. The tension is at \ang{90} to the motion: $W = 0$.}
\end{formulacard}

\textbf{Total energy.} The \term{total mechanical energy} is $E = K + U$. In a horizontal circle
the height does not change, so $U$ is constant. No net work is done, so $K$ is constant. So $E$
is constant too. A satellite in a circular orbit is the same. Gravity is at \ang{90} to its
motion, and the radius does not change (see \cref{ch5:chap}).

(In real life, air resistance slowly takes energy from a whirled ball. To keep the speed steady,
your hand moves in a small circle a little ahead of the ball. Then the string is not exactly at
\ang{90} to the motion. It pulls a little forwards and puts the lost energy back.)

\textbf{How work changes the total energy.} The work done by gravity is already counted in $U$.
Any \emph{other} force that does work changes the total energy: $W_\mathrm{other} = \Delta E$.
In a horizontal circle or a circular orbit, $W_\mathrm{other} = 0$, so $E$ stays constant.

A Ferris wheel is different. It turns at a steady speed in a vertical circle. So $K$ stays the
same, but $U$ grows as a rider rises. The seat pushes the rider up and does positive work. This
energy comes from the motor.

Example: a \qty{60}{kg} rider on a wheel of radius \qty{20}{m} goes from the bottom to the top.
The rider rises by the diameter, $2r = \qty{40}{m}$. So
$\Delta U = mg\Delta h = 60 \times 9.8 \times 40 = \qty{2.4e4}{J}$. The seat does
\qty{2.4e4}{J} of work on the rider, and gravity does $-\qty{2.4e4}{J}$. The net work is zero,
so $\Delta K = 0$.

\begin{formulacard}[notonsheet]{Work and total energy}
  \cardformula{W_\mathrm{other} = \Delta E = \Delta K + \Delta U}
  \begin{symbols}
    \sym{$W_\mathrm{other}$}{work done by all forces except gravity}{\si{J}}
    \sym{$\Delta E$}{change in total mechanical energy}{\si{J}}
    \sym{$\Delta K$}{change in kinetic energy}{\si{J}}
    \sym{$\Delta U$}{change in gravitational potential energy}{\si{J}}
  \end{symbols}
  \inwords{forces other than gravity (a push, a motor, friction, air resistance) change the
    total energy by the work they do.}
  \whentouse{to decide whether $E = K + U$ stays constant, or to find the work done by a motor,
    a push or friction.}
  \whennottouse{to find a force at one instant. Use $F_\mathrm{c} = \dfrac{mv^2}{r}$ for that.}
  \rearranged{$\Delta K = W_\mathrm{other} - \Delta U$; \ if $W_\mathrm{other} = 0$, $K + U$ stays constant}
  \tensecond{air resistance does $-\qty{0.30}{J}$ of work on a whirled ball. Its total energy
    falls by \qty{0.30}{J}.}
  \derive{start from the work--energy theorem.
    \begin{align*}
      W_\text{by gravity} + W_\mathrm{other} &= \Delta K
        && \text{net work $=$ gravity's work $+$ other work}\\
      -\Delta U + W_\mathrm{other} &= \Delta K
        && \text{gravity's work is the loss of $U$: $W_\text{by gravity} = -\Delta U$}\\
      W_\mathrm{other} &= \Delta K + \Delta U = \Delta E
        && \text{add $\Delta U$ to both sides}
    \end{align*}}
  \trap{do not count gravity twice. Its work is already in $\Delta U$, so leave it out of
    $W_\mathrm{other}$.}
\end{formulacard}

\begin{workedexample}{Speeding up on a roundabout}\label{ch3:we:roundabout}
A \qty{1200}{kg} car drives around a flat roundabout of radius \qty{25}{m}. (a) At a steady
\qty{18}{km.h^{-1}}, find the friction force towards the centre and the work it does in one
lap. (b) The driver then speeds up to \qty{36}{km.h^{-1}} during half a lap. Find the net work
done on the car and the average tangential (forward) force. Ignore air resistance.
\given{$m = \qty{1200}{kg}$, \ $r = \qty{25}{m}$, \ $u = 18 \div 3.6 = \qty{5.0}{m.s^{-1}}$, \
  $v = 36 \div 3.6 = \qty{10.0}{m.s^{-1}}$; \ half a lap: distance along the path
  $s = \pi r = \qty{78.54}{m}$ (not the straight-line displacement $2r = \qty{50}{m}$)}
\find{(a) $F_\mathrm{c}$ and its work \quad (b) $W_\mathrm{net}$ and the tangential force $F_\parallel$}
\formula{$F_\mathrm{c} = \dfrac{mv^2}{r}$, \quad $W_\mathrm{net} = \Delta K = \tfrac{1}{2}mv^2 - \tfrac{1}{2}mu^2$,
  \quad $W = F_\parallel s$}
\substitute{\eqlines{\text{(a) } F_\mathrm{c} &= \qty{1200}{kg} \times (\qty{5.0}{m.s^{-1}})^2 \div \qty{25}{m}
    = \qty{1200}{N}\\
  \text{work by } F_\mathrm{c} &= 0 \qquad \text{(the force is at \ang{90} to the motion)}\\
  \text{(b) } W_\mathrm{net} &= \tfrac{1}{2}(\qty{1200}{kg})(\qty{10.0}{m.s^{-1}})^2
    - \tfrac{1}{2}(\qty{1200}{kg})(\qty{5.0}{m.s^{-1}})^2\\
  &= \qty{6.0e4}{J} - \qty{1.5e4}{J} = \qty{4.5e4}{J}\\
  F_\parallel &= W_\mathrm{net} \div s = \qty{4.5e4}{J} \div \qty{78.54}{m} = \qty{573.0}{N}}\par
  (We rearranged $W = F_\parallel s$. $F_\parallel$ always points along the path, so we add up
  the tiny steps: $s$ is the distance along the arc.)}
\answer{(a) \final{$F_\mathrm{c} = \qty{1.2e3}{N}$}, \ work $= \qty{0}{J}$ \quad
  (b) \final{$W_\mathrm{net} = \qty{4.5e4}{J}$}, \ \final{$F_\parallel = \qty{5.7e2}{N}$}}
\checkstep{the forward acceleration is $a = F_\parallel \div m = 573.0 \div 1200 =
  \qty{0.4775}{m.s^{-2}}$. Then $v^2 = u^2 + 2as = 5.0^2 + 2 \times 0.4775 \times 78.54 = 100$, so
  $v = \qty{10.0}{m.s^{-1}}$, as given \cmark. At the end, $a_\mathrm{c} = 10.0^2 \div 25 =
  \qty{4.0}{m.s^{-2}}$ (0.41\,$g$), which a dry road can supply. While the car speeds up, its
  motion is not uniform.}
\end{workedexample}

\subsection{Vertical circles}

Swing a ball on a string in a vertical circle, or ride a rollercoaster loop. If only gravity
does work (no motor, no push from your hand, no friction), the speed is slowest at the top and
fastest at the bottom. This is \term{non-uniform circular motion} (moving in a circle with a
changing speed).

\begin{sidebyside}[0.5]
  \centering
  \begin{tikzpicture}[x=1cm, y=1cm]
    \def\R{1.7}
    \draw[hscmuted!60, line width=0.8pt] (0,0) circle (\R);
    \draw[hscink!60, line width=0.6pt] (0,0) -- (0,\R);
    \draw[hscink!60, line width=0.6pt] (0,0) -- (\R,0);
    \draw[hscink!60, line width=0.6pt] (0,0) -- (0,-\R);
    \fill[hscink] (0,0) circle (1.5pt);
    \node[diagramlabel, left, inner sep=3pt] at (0,0) {centre};
    \draw[->, hscmuted, line width=0.6pt] (155:2.0) arc (155:205:2.0);
    \node[diagramlabel, text=hscmuted, anchor=east] at (180:2.1) {motion};
    % top
    \draw[force] (-0.11,\R) -- (-0.11,1.1) node[left, inner sep=1pt, text=hscforce] {$F_\mathrm{T}$};
    \draw[force] (0.11,\R) -- (0.11,0.85) node[right, inner sep=1pt, text=hscforce] {$mg$};
    \draw[velocity] (0,\R) -- (-0.7,\R) node[left, inner sep=2pt, text=hscvel] {$v$};
    \fill[object] (0,\R) circle (4.5pt);
    \node[diagramlabel, above, inner sep=6pt] at (0,\R) {slowest};
    % side
    \draw[force] (\R,0) -- (0.9,0) node[above, inner sep=2pt, text=hscforce] {$F_\mathrm{T}$};
    \draw[force] (\R,0) -- (\R,-0.85) node[right, inner sep=2pt, text=hscforce] {$mg$};
    \draw[velocity] (\R,0) -- (\R,1.05) node[right, inner sep=2pt, text=hscvel] {$v$};
    \fill[object] (\R,0) circle (4.5pt);
    % bottom
    \draw[force] (0,-\R) -- (0,-0.4) node[right, pos=0.6, inner sep=2pt, text=hscforce] {$F_\mathrm{T}$};
    \draw[force] (0,-\R) -- (0,-2.55) node[right, inner sep=2pt, text=hscforce] {$mg$};
    \draw[velocity] (0,-\R) -- (1.5,-\R) node[right, inner sep=2pt, text=hscvel] {$v$};
    \fill[object] (0,-\R) circle (4.5pt);
    \node[diagramlabel, anchor=north east, inner sep=3pt] at (-0.15,-\R) {fastest};
  \end{tikzpicture}
  \captionof{figure}{A ball on a string in a vertical circle. Arrows are not to scale: at the
    bottom the tension is much bigger than $mg$.}\label{ch3:fig:vertical}
\nextside
  Look at \cref{ch3:fig:vertical}.
  \begin{itemize}
    \item The tension $F_\mathrm{T}$ (the pull of the string, \cref{ch2:chap}) always points
          along the string, to the centre. It is at \ang{90} to the motion, so it does no work.
    \item Gravity always points down. At the top and the bottom it is at \ang{90} to the motion.
    \item Everywhere else, part of the weight is along the path. Going up, it points backwards
          and slows the ball. Going down, it points forwards and speeds the ball up.
  \end{itemize}
\end{sidebyside}

So for a ball on a string (with no push from your hand), only gravity does work. This work
equals the loss of potential energy.
With no friction or air resistance, energy moves between $K$ and $U$, but the total stays the
same. This is the \term{conservation of mechanical energy}. From the top of a circle to the
bottom, the object falls through the diameter, $2r$.

\begin{formulacard}[combine={$K = \frac{1}{2}mv^2$ and $\Delta U = mg\Delta h$}]{Energy along a vertical circle}
  \cardformula{\tfrac{1}{2}mv_1^2 + mgh_1 = \tfrac{1}{2}mv_2^2 + mgh_2}
  \begin{symbols}
    \sym{$m$}{mass of the object}{\si{kg}}
    \sym{$v_1$, $v_2$}{speed at point 1 and at point 2}{\si{m.s^{-1}}}
    \sym{$h_1$, $h_2$}{height of point 1 and of point 2}{\si{m}}
    \sym{$g$}{gravitational acceleration, \qty{9.8}{m.s^{-2}}}{\si{m.s^{-2}}}
  \end{symbols}
  \inwords{kinetic energy plus potential energy is the same at every point of the path.}
  \whentouse{speeds at different heights on a vertical circle, a rollercoaster or a pendulum,
    when friction and air resistance are small.}
  \whennottouse{friction, air resistance or a motor does a lot of work. Also, it gives speeds,
    not forces.}
  \rearranged{$v_2 = \sqrt{v_1^2 + 2g(h_1 - h_2)}$; \ top to bottom of a circle:
    $v_\mathrm{bottom}^2 = v_\mathrm{top}^2 + 4gr$}
  \tensecond{$v_\mathrm{top} = \qty{3.0}{m.s^{-1}}$, $r = \qty{0.50}{m}$:
    $v_\mathrm{bottom} = \sqrt{3.0^2 + 4 \times 9.8 \times 0.50} = \qty{5.3}{m.s^{-1}}$.}
  \derive{the tension (or normal force) is at \ang{90} to the motion, so it does no work. With no
    friction, $W_\mathrm{other} = 0$, so $\Delta E = 0$ (``Work and total energy'' card).
    \begin{align*}
      \tfrac{1}{2}mv_1^2 + mgh_1 &= \tfrac{1}{2}mv_2^2 + mgh_2
        && \text{$E$ is the same at points 1 and 2}\\
      \tfrac{1}{2}mv_2^2 &= \tfrac{1}{2}mv_1^2 + mg(h_1 - h_2)
        && \text{subtract $mgh_2$ from both sides}\\
      v_2^2 &= v_1^2 + 2g(h_1 - h_2)
        && \text{multiply both sides by $\dfrac{2}{m}$}\\
      v_\mathrm{bottom}^2 &= v_\mathrm{top}^2 + 2g(2r) = v_\mathrm{top}^2 + 4gr
        && \text{top to bottom: $h_1 - h_2 = 2r$}
    \end{align*}}
  \trap{from the top to the bottom the height falls by $2r$, not $r$. The mass cancels, so you
    do not need $m$ to find a speed.}
\end{formulacard}

\textbf{Forces at the top and bottom.} To find a force, first find the speed at that point.

The net force towards the centre must be $F_\mathrm{c} = \dfrac{mv^2}{r}$ (on the sheet). It is
not an extra force. It is the sum of the real forces.

At the top and the bottom, every force is vertical. Take ``towards the centre'' as positive. A
force that points to the centre is $+$. A force that points away from the centre is $-$.
\Cref{ch3:fig:track} shows a cart in a dip and on a crest.

\begin{figure}
  \centering
  \begin{tikzpicture}[x=1cm, y=1cm]
    \draw[guide] (-0.2,0) -- (10.4,0);
    \node[diagramlabel, anchor=south west, text=hscmuted, font=\footnotesize] at (-0.2,0) {ground level};
    % track (heights to scale: 1 cm = 10 m)
    \draw[hscink, line width=1.4pt] (0,3.2) .. controls (1.5,3.2) and (2.0,0) .. (3.5,0)
      .. controls (5.0,0) and (5.6,2.0) .. (7.2,2.0)
      .. controls (8.6,2.0) and (9.2,0.6) .. (10,0.6);
    % A
    \fill[object] (0.2,3.32) circle (3.5pt);
    \node[diagramlabel, anchor=south west, inner sep=4pt, align=left] at (0.2,3.32) {\textbf{A}: at rest, \qty{32}{m} high};
    % B: dip
    \coordinate (B) at (3.5,0.12);
    \fill[hscink] (3.5,1.85) circle (1.5pt);
    \node[diagramlabel, anchor=west, font=\footnotesize, align=left] at (3.6,1.85) {centre of the dip\\(above B)};
    \draw[force] (B) -- (3.5,1.4) node[left, pos=0.7, inner sep=2pt, text=hscforce] {$F_\mathrm{N}$};
    \draw[force] (B) -- (3.5,-0.45) node[left, inner sep=2pt, text=hscforce] {$mg$};
    \fill[object] (B) circle (3.5pt);
    \node[diagramlabel, anchor=north west, align=left, font=\footnotesize] at (3.75,-0.08) {\textbf{B}: dip, $r = \qty{25}{m}$;\\net force up};
    % C: crest
    \coordinate (C) at (7.2,2.12);
    \fill[hscink] (7.2,1.05) circle (1.5pt);
    \node[diagramlabel, anchor=north, font=\footnotesize, align=center, inner sep=2pt] at (7.2,0.98) {centre of the crest\\(below C)};
    \draw[force] (C) -- (7.2,2.45) node[left, inner sep=2pt, text=hscforce] {$F_\mathrm{N}$};
    \draw[force] (C) -- (7.2,1.3) node[right, inner sep=2pt, text=hscforce] {$mg$};
    \fill[object] (C) circle (3.5pt);
    \node[diagramlabel, anchor=south west, align=left, font=\footnotesize] at (7.35,2.2) {\textbf{C}: crest, \qty{20}{m} high,\\$r = \qty{30}{m}$; net force down};
  \end{tikzpicture}
  \caption{A cart in a dip (B) and on a crest (C). The centre of the circle is above a dip and
    below a crest, so the net force points up at B and down at C. (Heights to scale; radii and
    arrows not to scale.)}\label{ch3:fig:track}
\end{figure}

\begin{itemize}
  \item String, top: the centre is below the ball. $F_\mathrm{T}$ and $mg$ both point down, to
        the centre. Both are $+$.
  \item String, bottom: the centre is above the ball. $F_\mathrm{T}$ points up ($+$). $mg$ points
        down, away ($-$).
  \item Cart in a \term{dip} (the low point of a valley in the track): the centre is above the
        cart. The \term{normal force} $F_\mathrm{N}$ (the push of the track or seat) is $+$, and
        $mg$ is $-$.
  \item Cart on a \term{crest} (the top of a hill in the track): the centre is below the cart.
        $mg$ is $+$, and $F_\mathrm{N}$ is $-$.
\end{itemize}


\begin{formulacard}[notonsheet]{Forces at the top and bottom of a vertical circle}
  \cardformula{\text{top: } F_\mathrm{T} + mg = \frac{mv^2}{r} \qquad
    \text{bottom: } F_\mathrm{T} - mg = \frac{mv^2}{r}}
  \begin{symbols}
    \sym{$F_\mathrm{T}$}{tension in the string (use $F_\mathrm{N}$ for a track or seat)}{\si{N}}
    \sym{$m$}{mass of the object}{\si{kg}}
    \sym{$g$}{gravitational acceleration, 9.8}{\si{m.s^{-2}}}
    \sym{$v$}{speed at that point}{\si{m.s^{-1}}}
    \sym{$r$}{radius of the circle}{\si{m}}
  \end{symbols}
  \inwords{forces towards the centre minus forces away from it $= mv^2/r$.}
  \whentouse{only at the very top or bottom, where every force is vertical.}
  \whennottouse{at the side of the circle, or to find a speed (use energy).}
  \rearranged{top: $F_\mathrm{T} = \dfrac{mv^2}{r} - mg$; \ bottom: $F_\mathrm{T} = mg + \dfrac{mv^2}{r}$;
    \ dip: $F_\mathrm{N} = mg + \dfrac{mv^2}{r}$; \ crest: $F_\mathrm{N} = mg - \dfrac{mv^2}{r}$;
    \ smallest speed at the top: $v = \sqrt{gr}$}
  \tensecond{a \qty{0.20}{kg} ball at the bottom, $r = \qty{0.50}{m}$, $v = \qty{4.0}{m.s^{-1}}$:
    $F_\mathrm{T} = 0.20 \times 9.8 + \dfrac{0.20 \times 4.0^2}{0.50} = 1.96 + 6.4 = \qty{8.4}{N}$.}
  \derive{$F_\mathrm{net} = F_\mathrm{c} = \dfrac{mv^2}{r}$ (on the sheet). Add the real forces,
    with ``towards the centre'' $+$. Top: $(+F_\mathrm{T}) + (+mg)$. Bottom: $(+F_\mathrm{T}) + (-mg)$.}
  \trap{at the top, $mg$ points \emph{to} the centre, so it is $+mg$, not $-mg$.}
\end{formulacard}

At the top, the string goes \emph{slack} (loose, with no pull) when $F_\mathrm{T} = 0$. Then
gravity alone gives the centripetal force:
\begin{align*}
  mg &= \frac{mv^2}{r} && \text{top, with } F_\mathrm{T} = 0\\
  g &= \frac{v^2}{r} && \text{divide both sides by } m\\
  v^2 &= gr && \text{multiply both sides by } r\\
  v &= \sqrt{gr} && \text{take the square root}
\end{align*}
This is the smallest speed at the top. Water in a swung bucket stays in for the same reason
when $v \geq \sqrt{gr}$. A cart on a crest uses the same steps: set $F_\mathrm{N} = 0$ in
$mg - F_\mathrm{N} = \dfrac{mv^2}{r}$. So a cart leaves a crest if its speed is more than
$\sqrt{gr}$.

On a ride you feel the seat's push $F_\mathrm{N}$. This is your \term{apparent weight}. You feel
heavier in a dip ($F_\mathrm{N} > mg$) and lighter on a crest ($F_\mathrm{N} < mg$).

\begin{workedexample}{Is the rollercoaster safe?}\label{ch3:we:coaster}
A rollercoaster cart with riders has a mass of \qty{650}{kg}. It starts from rest at A,
\qty{32}{m} above the ground. It rolls down into a dip B at ground level (radius \qty{25}{m}).
Then it goes over a crest C that is \qty{20}{m} high (radius \qty{30}{m}); see \cref{ch3:fig:track}.
Assume no friction or air resistance. (a) Find the normal force on the cart at B. (b) Show
that the cart stays on the track at C, and find the normal force there.
\given{$m = \qty{650}{kg}$, \ $v_\mathrm{A} = 0$, \ $h_\mathrm{A} = \qty{32}{m}$; \
  $h_\mathrm{B} = 0$, $r_\mathrm{B} = \qty{25}{m}$; \ $h_\mathrm{C} = \qty{20}{m}$, $r_\mathrm{C} = \qty{30}{m}$}
\find{(a) $F_\mathrm{N}$ at B \quad (b) whether $F_\mathrm{N} \geq 0$ at C, and its value}
\formula{energy from A: $v^2 = 2g(h_\mathrm{A} - h)$ (from $\tfrac12 mv_\mathrm{A}^2 + mgh_\mathrm{A} =
  \tfrac12 mv^2 + mgh$ with $v_\mathrm{A} = 0$: divide by $m$, multiply by 2, then subtract $2gh$); \
  dip: $F_\mathrm{N} = mg + \dfrac{mv^2}{r}$; \ crest: $F_\mathrm{N} = mg - \dfrac{mv^2}{r}$}
\substitute{\eqlines{v_\mathrm{B}^2 &= 2(\qty{9.8}{m.s^{-2}})(\qty{32}{m} - \qty{0}{m}) = \qty{627.2}{m^2.s^{-2}}
    \qquad \text{(A to B)}\\
  F_\mathrm{N,B} &= mg + \frac{mv_\mathrm{B}^2}{r_\mathrm{B}} \qquad \text{(dip)}\\
  &= (\qty{650}{kg})(\qty{9.8}{m.s^{-2}}) + \frac{(\qty{650}{kg})(\qty{627.2}{m^2.s^{-2}})}{\qty{25}{m}}\\
  &= \qty{6370}{N} + \qty{16307}{N} = \qty{22677}{N}}\par
  \eqlines{v_\mathrm{C}^2 &= 2(\qty{9.8}{m.s^{-2}})(\qty{32}{m} - \qty{20}{m}) = \qty{235.2}{m^2.s^{-2}}
    \qquad \text{(A to C)}\\
  v_\mathrm{C} &= \sqrt{\qty{235.2}{m^2.s^{-2}}} = \qty{15.34}{m.s^{-1}}\\
  \sqrt{g r_\mathrm{C}} &= \sqrt{(\qty{9.8}{m.s^{-2}})(\qty{30}{m})} = \qty{17.15}{m.s^{-1}}
    \qquad \text{(speed limit at C)}}\par
  \eqlines{F_\mathrm{N,C} &= mg - \frac{mv_\mathrm{C}^2}{r_\mathrm{C}} \qquad \text{(crest)}\\
  &= \qty{6370}{N} - \frac{(\qty{650}{kg})(\qty{235.2}{m^2.s^{-2}})}{\qty{30}{m}}\\
  &= \qty{6370}{N} - \qty{5096}{N} = \qty{1274}{N}}}
\answer{(a) \final{$F_\mathrm{N} = \qty{2.3e4}{N}$} at B (3.6 times the weight) \quad
  (b) $v_\mathrm{C} < \sqrt{gr_\mathrm{C}}$, so it stays on: \final{$F_\mathrm{N} = \qty{1.3e3}{N}$}
  at C (0.20 times the weight)}
\checkstep{$F_\mathrm{N} > mg$ in the dip (riders feel heavy) and $F_\mathrm{N} < mg$ on the crest
  (riders feel light) \cmark. A crest radius below $\dfrac{v_\mathrm{C}^2}{g} = \qty{24}{m}$ would
  throw the cart off.}
\end{workedexample}

\begin{mistakes}
\begin{itemize}
  \item \xmark\ ``The centripetal force does work: force $\times$ distance round the circle.''
        \cmark\ It is at \ang{90} to the motion, so $W = 0$.
  \item \xmark\ Height change $r$ from the top to the bottom of a circle. \cmark\ It is $2r$.
  \item \xmark\ $F_\mathrm{T} - mg = \dfrac{mv^2}{r}$ at the top. \cmark\ Both point to the centre:
        $F_\mathrm{T} + mg$.
  \item \xmark\ ``$K + U$ is constant in every vertical circle.'' \cmark\ Only if no other force
        does work. A Ferris wheel's motor adds energy.
\end{itemize}
\end{mistakes}

\begin{examtip}
``Explain why the kinetic energy in uniform circular motion is constant'' (3 marks). Make three
linked points. (1) The net force points to the centre, at \ang{90} to the velocity at every
instant. (2) So $W = Fs\cos\ang{90} = 0$. (3) By the work--energy theorem, $\Delta K = 0$, so the
speed is constant. If the question asks about \emph{total} energy, add (4): the height (or orbit
radius) does not change, so $U$ is constant, and so $E = K + U$ is constant.
Vertical-circle questions have two steps. First use energy to find the speed (write ``no
friction''). Then use the forces at the top or bottom. End with a comment: heavier or lighter,
stays on or not.
\end{examtip}

\begin{quickcheck}
  \item A \qty{0.050}{kg} puck slides on a smooth horizontal table in a circle of radius
        \qty{0.60}{m} at \qty{5.0}{m.s^{-1}}. A string ties it to a pin at the centre. Find (a) the
        tension, (b) the work done by the tension in 10 turns and (c) the puck's kinetic energy.
  \item Explain why a satellite in a circular orbit keeps a constant speed and a constant total energy.
  \item A bucket of water is swung in a vertical circle of radius \qty{1.1}{m}. (a) Find the
        smallest speed at the top that keeps the water in. (b) Find the speed at the bottom if the
        top speed has this value.
  \item A Ferris wheel of radius \qty{15}{m} turns at a steady speed. A \qty{50}{kg} rider goes from
        the bottom to the top. (a) Does the rider's kinetic energy change? (b) Find the change in
        potential energy. (c) Which force does the work that supplies this energy?
\end{quickcheck}

%% ==================================================================
\section{Torque: what makes things turn}\label{ch3:sec:torque}
\covers{M5.IQ2.05}

\begin{bigidea}
A force can make an object turn about a pivot. Its turning effect is the torque,
$\tau = r_\perp F$. A big force, far from the pivot and at \ang{90} to the lever, gives a big
torque. If the torques balance, the rotation does not change.
\end{bigidea}

\begin{keywords}
  \kw{Torque ($\tau$)}{the turning effect of a force. Unit: newton metre (\si{N.m}).}
  \kw{Pivot (axis)}{the point or line that an object turns about, such as a hinge.}
  \kw{Lever}{a bar that turns about a pivot, such as a spanner or a door.}
  \kw{Line of action}{the line along which a force acts, drawn longer in both directions.}
  \kw{Perpendicular distance ($r_\perp$)}{the shortest distance from the pivot to the line of
    action of the force.}
\end{keywords}

Push a door next to its hinges and it hardly moves. Push at the handle, far from the hinges,
and it swings easily. Push straight towards the hinges and it does not turn at all. The turning
effect of a force is its \term{torque} ($\tau$, ``tau''). The point it turns about is the
\term{pivot}.

Torque is bigger when:
\begin{itemize}
  \item the force is bigger;
  \item the force acts further from the pivot;
  \item the force is closer to \ang{90} to the \term{lever} (the bar that turns).
\end{itemize}

\begin{toolkit}{Sine in torque problems}
Splitting a force (toolkit in \cref{ch1:sec:model}): with $\theta$ measured from the lever, the
part of $F$ at \ang{90} to the lever is opposite $\theta$, so it is $F\sin\theta$ (SOH). The part
along the lever is next to $\theta$, so it is $F\cos\theta$ (CAH).
\begin{itemize}
  \item $\sin\ang{90} = 1$ gives the most torque. $\sin\ang{0} = \sin\ang{180} = 0$, so a force
        along the lever gives none.
  \item $\sin\theta = \sin(\ang{180} - \theta)$, for example $\sin\ang{30} = \sin\ang{150} = 0.5$.
        So you may measure the angle on either side of the force: $\theta$ and
        $\ang{180} - \theta$ give the same torque.
  \item If the angle is measured from a line at \ang{90} to the lever, use \ang{90} minus that angle.
  \item The angles here are in degrees. Set your calculator to degree mode (DEG): $\sin\ang{90}$
        must give 1. In radian mode, $\sin 90$ gives 0.894, which is wrong here.
\end{itemize}
\end{toolkit}

\Cref{ch3:fig:torque} shows two ways to find a torque. Both give the same answer. In both, $r$
goes from the pivot to the point where the force acts. $\theta$ is the angle between $r$ and
the force.
\begin{itemize}
  \item \textbf{(a) Perpendicular distance.} Extend the force into a line: its \term{line of
        action}. Draw a line from the pivot that meets it at \ang{90}. This is the
        \term{perpendicular distance}, $r_\perp = r\sin\theta$. Then $\tau = r_\perp F$.
  \item \textbf{(b) Perpendicular force.} Split $F$ into $F\sin\theta$ (at \ang{90} to the lever)
        and $F\cos\theta$ (along the lever). Only $F\sin\theta$ turns, so $\tau = r \times F\sin\theta$.
\end{itemize}

\begin{figure}
  \centering
  \begin{tikzpicture}[x=1cm, y=1cm]
    % ---- (a) perpendicular distance
    \begin{scope}
      \draw[object, rounded corners=1.5pt] (0.25,-0.13) rectangle (3.25,0.13);
      \draw[object] (0:0.34) -- (60:0.34) -- (120:0.34) -- (180:0.34) -- (240:0.34) -- (300:0.34) -- cycle;
      \fill[hscink] (0,0) circle (1.6pt);
      \node[diagramlabel, below, inner sep=5pt] at (0,-0.3) {pivot};
      \coordinate (P) at (3.0,0);
      \coordinate (Q) at (2.25,-1.299);
      \draw[guide] (P) -- (3.9,0);
      \draw[guide] ($(Q)+(240:0.35)$) -- (P);
      \draw[hscbrand, line width=1.1pt] (0,0) -- node[below left, diagramlabel, text=hscbrand, inner sep=1pt] {$r_\perp$} (Q);
      \draw[hscink, line width=0.5pt] ($(Q)+(150:0.2)$) -- ($(Q)+(150:0.2)+(60:0.2)$) -- ($(Q)+(60:0.2)$);
      \draw[hscink, <->, line width=0.5pt] (0,0.42) -- node[above, diagramlabel] {$r$} (3.0,0.42);
      \draw[hscink, line width=0.6pt] ($(P)+(0:0.5)$) arc (0:60:0.5);
      \node[diagramlabel] at ($(P)+(30:0.75)$) {$\theta$};
      \draw[force] (P) -- ++(60:1.3) node[above, inner sep=2pt, text=hscforce] {$F$};
      \fill[hscink] (P) circle (1.4pt);
      \node[diagramlabel] at (1.6,-1.95) {\textbf{(a)} $r_\perp = r\sin\theta$, \ $\tau = r_\perp F$};
    \end{scope}
    % ---- (b) perpendicular force
    \begin{scope}[xshift=6.9cm]
      \draw[object, rounded corners=1.5pt] (0.25,-0.13) rectangle (3.25,0.13);
      \draw[object] (0:0.34) -- (60:0.34) -- (120:0.34) -- (180:0.34) -- (240:0.34) -- (300:0.34) -- cycle;
      \fill[hscink] (0,0) circle (1.6pt);
      \node[diagramlabel, below, inner sep=5pt] at (0,-0.3) {pivot};
      \draw[hscink, <->, line width=0.5pt] (0,0.42) -- node[above, diagramlabel] {$r$} (2.92,0.42);
      \coordinate (P) at (3.0,0);
      \draw[guide] ($(P)+(0.65,0)$) -- ($(P)+(60:1.3)$) -- ($(P)+(0,1.126)$);
      \draw[force, hscforce!40, line width=1pt] (P) -- ++(60:1.3) node[right, inner sep=2pt, text=hscforce!70] {$F$};
      \draw[force] (P) -- ++(0,1.126) node[left, inner sep=5pt, text=hscforce, align=right] {$F\sin\theta$\\(turns)};
      \draw[force] (P) -- ++(0.65,0) node[right, inner sep=3pt, text=hscforce, align=left] {$F\cos\theta$\\(no turning)};
      \draw[hscink, line width=0.6pt] ($(P)+(0:0.35)$) arc (0:60:0.35);
      \node[diagramlabel] at ($(P)+(25:0.58)$) {$\theta$};
      \fill[hscink] (P) circle (1.4pt);
      \node[diagramlabel] at (1.6,-1.95) {\textbf{(b)} $\tau = r \times F\sin\theta$};
    \end{scope}
  \end{tikzpicture}
  \caption{A force $F$ on a spanner at angle $\theta$ to the handle. (a) The perpendicular
    distance from the pivot to the line of action. (b) Only the part of $F$ at \ang{90} to the
    handle turns the nut.}\label{ch3:fig:torque}
\end{figure}

\begin{formulacard}{Torque}
  \cardformula{\tau = r_\perp F = rF\sin\theta}
  \begin{symbols}
    \sym{$\tau$}{torque: turning effect of the force}{\si{N.m}}
    \sym{$r_\perp$}{perpendicular distance from the pivot to the line of action}{\si{m}}
    \sym{$F$}{size of the force}{\si{N}}
    \sym{$r$}{distance from the pivot to where the force acts}{\si{m}}
    \sym{$\theta$}{angle between $r$ and the force}{\si{\degree} (calculator in degree mode)}
  \end{symbols}
  \inwords{turning effect = force $\times$ perpendicular distance from the pivot.}
  \whentouse{a force turns an object about a pivot: doors, spanners, pedals, seesaws, motor
    coils (Module 6).}
  \whennottouse{the angle you are given is measured from a line at \ang{90} to the lever.
    Change it to the angle between $r$ and $F$ first ($\theta = \ang{90} - $ that angle).}
  \rearranged{$F = \dfrac{\tau}{r\sin\theta}$ \qquad $r = \dfrac{\tau}{F\sin\theta}$ \qquad
    $\sin\theta = \dfrac{\tau}{rF}$}
  \tensecond{\qty{50}{N} at \ang{90} to the end of a \qty{0.30}{m} spanner:
    $\tau = 0.30 \times 50 \times \sin\ang{90} = \qty{15}{N.m}$.}
  \trap{measure $r$ from the pivot, not from the end of the object. The unit is \si{N.m}
    (not J, and not ``nm'').}
\end{formulacard}

\textbf{How torque changes rotation.} A force changes straight-line motion. In the same way, a
torque changes turning motion.
\begin{itemize}
  \item If the net torque is not zero, the object starts to turn, speeds up or slows down. A bigger
        net torque changes the rotation faster. A bicycle chain's torque speeds up the wheel. The
        brakes' opposite torque slows it down.
  \item If the torques are \term{balanced} (net torque zero), the rotation does not change. This is
        Newton's first law for turning: with no net torque, a still object stays still, and a
        turning object keeps turning at the same rate. The \term{clockwise} torques (turning the
        same way as a clock's hands) then equal the \term{anticlockwise} torques (turning the
        other way).
\end{itemize}

Take every torque about the same pivot.

Example: a \qty{20}{kg} child sits \qty{1.5}{m} from a seesaw's pivot. Her torque is
$20 \times 9.8 \times 1.5 = \qty{294}{N.m}$ (anticlockwise). A \qty{30}{kg} child sits
\qty{1.0}{m} from the pivot on the other side. His torque is $30 \times 9.8 \times 1.0 =
\qty{294}{N.m}$ (clockwise). The torques are equal, so the seesaw balances.

\begin{formulacard}[notonsheet]{Balanced torques}
  \cardformula{\sum\tau_\mathrm{clockwise} = \sum\tau_\mathrm{anticlockwise}}
  \begin{symbols}
    \sym{$\tau$}{each torque about the pivot (clockwise: the same way as a clock's hands;
      anticlockwise: the other way)}{\si{N.m}}
    \sym{$\sum$}{``add up all of them''}{no unit}
  \end{symbols}
  \inwords{if the turning effects each way are equal, the rotation does not change.}
  \whentouse{seesaws, signs on poles, a balanced metre rule: the rotation does not change.}
  \whennottouse{the rotation starts, speeds up or slows down (the net torque is not zero).}
  \rearranged{one force each side: $r_1F_1 = r_2F_2$, so $r_2 = \dfrac{r_1F_1}{F_2}$ and
    $F_2 = \dfrac{r_1F_1}{r_2}$}
  \tensecond{a \qty{250}{N} child \qty{2.0}{m} from the pivot is balanced by a \qty{400}{N}
    child at $r_2 = \dfrac{2.0 \times 250}{400} = \qty{1.25}{m}$ on the other side.}
  \derive{no change in rotation means the net torque is zero:
    $\sum\tau_\mathrm{clockwise} - \sum\tau_\mathrm{anticlockwise} = 0$. Add
    $\sum\tau_\mathrm{anticlockwise}$ to both sides.}
  \trap{take every torque about the \emph{same} pivot, and measure every $r$ from that pivot.}
\end{formulacard}

Everyday uses of torque:
\begin{itemize}
  \item door handles far from the hinges;
  \item long spanners, wide steering wheels and wide tap handles;
  \item bicycle pedals: the \term{crank} (the arm from the pedal to the axle) gives the most torque
        when it is horizontal and none when it is vertical, if the foot pushes straight down;
  \item electric motors (Module 6).
\end{itemize}

\begin{workedexample}{Turning a nut}\label{ch3:we:spanner}
A mechanic pushes with \qty{80}{N} on the end of a \qty{0.25}{m} spanner. Find the torque on the
nut when the force is (a) at \ang{90} to the spanner, (b) at \ang{60} to the spanner and (c) along
the spanner, towards the nut.
\given{$F = \qty{80}{N}$, \ $r = \qty{0.25}{m}$; \ $\theta = \ang{90}$, \ang{60} and \ang{180}
  (a push towards the nut points back along $r$)}
\find{the torque $\tau$ in each case}
\formula{$\tau = rF\sin\theta$ (calculator in degree mode)}
\substitute{\eqlines{\text{(a) } \tau &= \qty{0.25}{m} \times \qty{80}{N} \times \sin\ang{90} = 0.25 \times 80 \times 1 = \qty{20}{N.m}\\
  \text{(b) } \tau &= \qty{0.25}{m} \times \qty{80}{N} \times \sin\ang{60} = 0.25 \times 80 \times 0.866 = \qty{17.32}{N.m}\\
  \text{(c) } \tau &= \qty{0.25}{m} \times \qty{80}{N} \times \sin\ang{180} = \qty{0}{N.m}}}
\answer{(a) \final{$\tau = \qty{20}{N.m}$} \quad (b) \final{$\tau = \qty{17}{N.m}$} \quad
  (c) \final{$\tau = \qty{0}{N.m}$}}
\checkstep{with method (a) for part (b): $r_\perp = 0.25\sin\ang{60} = \qty{0.2165}{m}$, and
  $0.2165 \times 80 = \qty{17.3}{N.m}$ \cmark. The torque is largest at \ang{90}.}
\end{workedexample}

\begin{workedexample}{Holding up a plant basket}\label{ch3:we:sign}
\begin{sidebyside}[0.56]
A \qty{4.0}{kg} plant basket hangs from the end of a light (almost massless) horizontal rod,
\qty{1.5}{m} long. The rod is hinged to a wall. A cable joins the end of the rod to a point
higher up the wall, at \ang{35} to the rod (\cref{ch3:fig:sign}). Find the tension in the cable.
\nextside
  \centering
  \begin{tikzpicture}[x=0.9cm, y=0.9cm]
    \fill[ground] (-0.3,-1.9) rectangle (0,2.8);
    \draw[hscink, line width=0.7pt] (0,-1.9) -- (0,2.8);
    \draw[hscink, line width=0.8pt] (3.6,0) -- (0,2.52);
    \draw[hscink, line width=2.4pt] (0,0) -- (3.6,0);
    \draw[object] (0,0) circle (3pt);
    \draw[hscink, <->, line width=0.5pt] (0,-0.28) -- (3.6,-0.28);
    \node[diagramlabel, anchor=north, inner sep=2pt] at (1.5,-0.3) {$r = \qty{1.5}{m}$};
    \node[diagramlabel, anchor=north west, inner sep=2pt, font=\footnotesize] at (0.05,-0.95) {hinge (pivot)};
    \draw[hscink, line width=0.6pt] (3.6,0) -- (3.6,-0.42);
    \draw[hscink, line width=0.6pt] (3.6,-0.42) -- (3.05,-0.6);
    \draw[hscink, line width=0.6pt] (3.6,-0.42) -- (4.15,-0.6);
    \draw[object] (2.9,-0.6) -- (4.3,-0.6) -- (4.02,-1.2) -- (3.18,-1.2) -- cycle;
    \node[diagramlabel, font=\scriptsize] at (3.6,-0.9) {basket};
    \draw[force] (3.6,-1.2) -- (3.6,-1.95) node[left, inner sep=3pt, text=hscforce] {$mg$};
    \draw[force] (3.6,0) -- ++(145:1.3) node[above right, inner sep=1pt, text=hscforce] {$F_\mathrm{T}$};
    \draw[hscink, line width=0.6pt] ($(3.6,0)+(145:0.75)$) arc (145:180:0.75);
    \node[diagramlabel, font=\footnotesize] at ($(3.6,0)+(163:1.12)$) {\ang{35}};
    \node[diagramlabel, anchor=west] at (0.1,2.55) {cable};
  \end{tikzpicture}
  \captionof{figure}{A plant basket held up by a hinged rod and a cable.}\label{ch3:fig:sign}
\end{sidebyside}
\given{$m = \qty{4.0}{kg}$, so $mg = 4.0 \times 9.8 = \qty{39.2}{N}$; \ $r = \qty{1.5}{m}$ (hinge to
  the end); \ the weight acts at the end, at \ang{90} to the rod; \ the cable makes \ang{35} with
  the rod, so the angle between $r$ and $F_\mathrm{T}$ is \ang{145}}
\find{the tension $F_\mathrm{T}$ in the cable}
\formula{balanced torques about the hinge: $\tau_\mathrm{cable} = \tau_\mathrm{weight}$, with
  $\tau = rF\sin\theta$}
\substitute{\eqlines{\tau_\mathrm{weight} &= \qty{1.5}{m} \times \qty{39.2}{N} \times \sin\ang{90}
    = \qty{58.8}{N.m} \quad \text{(clockwise)}\\
  \tau_\mathrm{cable} &= \qty{1.5}{m} \times F_\mathrm{T} \times \sin\ang{145}
    \quad \text{(anticlockwise)}\\
  &= 1.5 \times 0.5736 \times F_\mathrm{T} = 0.8604\,F_\mathrm{T} \quad \text{($\sin\ang{145} = \sin\ang{35}$)}\\
  0.8604\,F_\mathrm{T} &= 58.8 \quad \text{(set the torques equal)}\\
  F_\mathrm{T} &= 58.8 \div 0.8604 = \qty{68.34}{N} \quad \text{(divide both sides by 0.8604)}}}
\answer{\final{$F_\mathrm{T} = \qty{68}{N}$} (2 significant figures)}
\checkstep{the hinge force has no torque about the hinge, so we chose it. About the rod's end,
  the weight and cable have no torque, so the hinge force has none either: it is horizontal.
  So the vertical part of the tension holds up the whole weight:
  $F_\mathrm{T}\sin\ang{35} = 68.34 \times 0.5736 = \qty{39.2}{N} = mg$ \cmark.}
\end{workedexample}

\begin{mistakes}
\begin{itemize}
  \item \xmark\ Using $\cos\theta$ when $\theta$ is between $r$ and $F$. \cmark\ Use $\sin\theta$.
        Draw the diagram and mark $\theta$ first.
  \item \xmark\ Measuring $r$ from the end of the object. \cmark\ Measure from the pivot.
  \item \xmark\ ``A bigger force always gives more torque.'' \cmark\ A force through the pivot gives
        zero torque, however big it is.
  \item \xmark\ Calculator in radian mode. \cmark\ Use degree mode for angles in degrees.
\end{itemize}
\end{mistakes}

\begin{examtip}
``Explain why door handles are far from the hinges'' (2--3 marks). Name the hinge as the pivot.
Quote $\tau = r_\perp F$. The door needs a certain torque, so a larger $r_\perp$ means a smaller
force. Pushing at \ang{90} ($\sin\ang{90} = 1$) gives the most torque. In calculations, draw the
pivot, $r$ and $\theta$. Substitute into $\tau = rF\sin\theta$ and give the unit \si{N.m}.
\end{examtip}

\begin{quickcheck}
  \item A \qty{120}{N} force acts on the end of a \qty{0.35}{m} spanner. Find the torque when the
        force is (a) at \ang{90} and (b) at \ang{70} to the spanner.
  \item A bicycle crank is \qty{0.17}{m} long. The rider pushes straight down on the pedal with
        \qty{600}{N}. Find the torque when the crank is (a) horizontal, (b) \ang{45} below the
        horizontal and (c) pointing straight down.
  \item A \qty{30}{kg} child sits \qty{1.8}{m} from the pivot of a seesaw. Where must a \qty{45}{kg}
        adult sit, on the other side, to balance it? Explain your reasoning.
\end{quickcheck}

%% ==================================================================
\section{Key experiments}\label{ch3:sec:experiments}
\covers{M5.IQ2.04, M5.IQ2.05}

The syllabus asks you to ``investigate'' both ideas in this chapter. Exams often ask you to
design, improve or judge tests like these.

\begin{experiment}{Investigating torque: $\tau = r_\perp F$}
\aim{To test $\tau = rF\sin\theta$. We find how the force that balances a fixed load depends on
  its distance from the pivot and on its angle. Then we see how an unbalanced torque changes the
  rotation.}
\apparatus{metre rule with a hole at the \qty{50}{cm} mark, pivot pin on a retort stand,
  \qty{200}{g} mass on a thread loop, slotted masses (\qty{50}{g}), spring balance
  (\qtyrange{0}{5}{N}), protractor (\cref{ch3:fig:torqueexp}).
  \begin{figure}
    \centering
    \begin{tikzpicture}[x=0.8cm, y=0.8cm]
      \fill[hscmuted!35] (-0.9,-2.45) rectangle (0.9,-2.25);
      \draw[hscmuted, line width=2.2pt] (0,-2.25) -- (0,0);
      \draw[object] (-4.5,-0.1) rectangle (4.5,0.1);
      \foreach \x in {-4.5,-3.6,...,4.6} \draw[hscink, line width=0.4pt] (\x,0.1) -- (\x,-0.02);
      \fill[hscink] (0,0) circle (2pt);
      \draw[hscink, line width=0.6pt] (-2.7,-0.1) -- (-2.7,-0.7);
      \draw[object] (-3.2,-0.7) rectangle (-2.2,-1.25);
      \node[diagramlabel, font=\scriptsize] at (-2.7,-0.975) {\qty{200}{g}};
      \draw[hscink, line width=0.6pt] (2.7,-0.1) -- (2.7,-0.4);
      \draw[object] (2.55,-0.4) rectangle (2.85,-1.5);
      \foreach \y in {-0.55,-0.7,...,-1.4} \draw[hscink, line width=0.3pt] (2.55,\y) -- (2.66,\y);
      \draw[force] (2.7,-1.5) -- (2.7,-2.2) node[right, inner sep=3pt, text=hscforce] {pull down with $F$};
      \draw[hscink, <->, line width=0.5pt] (-2.7,0.45) -- node[above, diagramlabel] {\qty{0.300}{m}} (0,0.45);
      \draw[hscink, <->, line width=0.5pt] (0,0.45) -- node[above, diagramlabel] {$r$} (2.7,0.45);
      \node[diagramlabel, anchor=south] at (0,1.15) {pivot at the \qty{50}{cm} mark};
      \node[diagramlabel, anchor=east] at (-3.3,-0.975) {load};
      \node[diagramlabel, anchor=west] at (2.95,-0.95) {spring balance};
    \end{tikzpicture}
    \caption{The metre rule balanced by a load and a spring balance.}\label{ch3:fig:torqueexp}
  \end{figure}}
\method{\begin{enumerate}[beginpenalty=10000]
  \item Hang the rule on the pivot at its \qty{50}{cm} mark. Check that the empty rule balances
        level. Then its \term{centre of mass} (the point where all its weight seems to act) is
        at the pivot.
        Hang the load \qty{0.300}{m} left of the pivot. Its torque is fixed:
        $0.300 \times 0.200 \times 9.8 = \qty{0.588}{N.m}$.
  \item Hook the spring balance $r = \qty{0.150}{m}$ right of the pivot. Pull straight down until
        the rule is level. Record $F$. Repeat three times and average.
  \item Repeat for $r$ from \qty{0.200}{m} to \qty{0.450}{m}. Then keep $r = \qty{0.400}{m}$ and
        pull at \ang{90}, \ang{60}, \ang{45} and \ang{30} to the rule.
  \item Rotation test: take the spring balance away. Hang \qty{150}{g} \qty{0.300}{m} right of the
        pivot. The torques no longer balance, so the rule starts to turn. Repeat with \qty{100}{g}
        and then \qty{50}{g}. Film each test to compare how fast the rule starts to turn.
\end{enumerate}
\textbf{Variables.} Independent (the one you change): $r$, then $\theta$. Dependent (the one you
measure): $F$. Controlled (kept the same): the load, its distance (\qty{0.300}{m}), the pivot,
and a level rule.}
\results{Sample results (each $F$ is the mean of three readings):
  \begin{sidebyside}[0.42]
  \centering\small
  \begin{tabular}{@{}cccc@{}}
    \toprule
    $r$ (\si{m}) & $1/r$ (\si{m^{-1}}) & $F$ (\si{N}) & $Fr$ (\si{N.m})\\
    \midrule
    0.150 & 6.67 & 3.95 & 0.593\\
    0.200 & 5.00 & 2.90 & 0.580\\
    0.250 & 4.00 & 2.35 & 0.588\\
    0.300 & 3.33 & 1.95 & 0.585\\
    0.350 & 2.86 & 1.70 & 0.595\\
    0.400 & 2.50 & 1.45 & 0.580\\
    0.450 & 2.22 & 1.30 & 0.585\\
    \bottomrule
  \end{tabular}
\nextside
  \centering
  \begin{tikzpicture}
    \begin{axis}[hsc, width=\linewidth, height=0.62\linewidth,
        xmin=0, xmax=7.5, ymin=0, ymax=4.5,
        xlabel={$1/r$ (\si{m^{-1}})}, ylabel={$F$ (\si{N})}]
      \addplot[hscseriesA, only marks, mark=*, mark size=2pt]
        coordinates {(6.667,3.95) (5.000,2.90) (4.000,2.35) (3.333,1.95) (2.857,1.70) (2.500,1.45) (2.222,1.30)};
      \addplot[hscseriesB, dashed, line width=1pt, domain=0:7.2, samples=2] {0.5876*x};
      \node[anchor=west, font=\footnotesize, text=hscink] at (axis cs:0.25,3.9)
        {gradient $= \qty{0.588}{N.m}$};
    \end{axis}
  \end{tikzpicture}
  \captionof{figure}{$F$ against $1/r$.}\label{ch3:fig:torquegraph}
\end{sidebyside}
  \textbf{Why plot $F$ against $1/r$?} The load's torque $\tau$ is fixed, and the balance must
  give the same torque:
  \begin{align*}
    Fr &= \tau && \text{torques balance}\\
    F &= \tau \times \frac{1}{r} && \text{divide both sides by } r
  \end{align*}
  This has the form $y = mx$, with $y = F$, $x = 1/r$ and gradient $m = \tau$. So the graph is a
  straight line through the origin, and its gradient is the torque.

  Gradient $=$ rise $\div$ run (toolkit in \cref{ch1:sec:model}). Use two points on the line,
  $(0, 0)$ and $(6.67, 3.92)$: gradient $= (3.92 - 0) \div (6.67 - 0) = \qty{0.588}{N.m}$, the
  same as the load's torque.

  At $r = \qty{0.400}{m}$ the angles \ang{90}, \ang{60}, \ang{45}, \ang{30} needed 1.45, 1.70, 2.10,
  \qty{2.95}{N}, so $F\sin\theta = 1.45$, 1.47, 1.48, \qty{1.48}{N}: the same within reading
  error (the small error in reading the scale).

  Rotation test: with \qty{150}{g}, \qty{100}{g} and \qty{50}{g} on the right, the net torque was
  0.147, 0.294 and \qty{0.441}{N.m}. The rule turned faster each time.}
\conclusion{$F \times r$ stays at about \qty{0.59}{N.m}, the load's torque. $F$ against $1/r$ is a
  straight line through the origin (gradient \qty{0.588}{N.m}). At an angle only $F\sin\theta$
  turns the rule. This supports $\tau = rF\sin\theta$. When the torques did not balance, the rule
  started to turn, and a bigger net torque changed its rotation faster.}
\significance{It explains door handles, spanners and levers, and motor torque (Module 6).}
\validity{\begin{itemize}[beginpenalty=10000]
  \item \textbf{Validity:} check that the empty rule balances level on the pivot, so its own
        weight has no torque. Measure $r$ from the pivot. Keep the rule level. Change one
        variable at a time.
  \item \textbf{Reliability:} take three readings at each distance and average them. The points
        lie close to one line.
  \item \textbf{Accuracy:} zero the spring balance in the position you use it. Read it at eye
        level. A force sensor is more accurate.
\end{itemize}}
\end{experiment}

\begin{experiment}{Work and energy in circular motion}
\aim{To investigate how the work done on an object moving in a circle is related to its total
  energy: (A) in a horizontal circle, where the tension is at \ang{90} to the motion; (B) in a
  vertical swing, where gravity does work.}
\apparatus{\textbf{Part A:} a \qty{0.025}{kg} rubber stopper on a string through a plastic tube,
  a \qty{0.100}{kg} hanging mass, a paper clip (marker), a phone camera that films at 240 frames
  per second, a metre rule. \textbf{Part B:} a \qty{0.0640}{kg} steel ball (diameter
  \qty{25.0}{mm}) on a \qty{1.0}{m} string, a clamp stand, a \term{light gate} (a light beam; the
  data logger times how long the ball blocks it) and data logger, a metre rule, vernier calipers
  (\cref{ch3:fig:pendulum}).
  \begin{figure}
    \centering
    \begin{tikzpicture}[x=0.75cm, y=0.75cm]
      % ---- (a) Part A: stopper in a horizontal circle
      \draw[displacement, dashed] (0,2.0) ellipse (2.1 and 0.5);
      \draw[hscink, line width=0.7pt] (0,2.0) -- (2.1,2.0);
      \node[diagramlabel, below, inner sep=2pt] at (1.05,2.0) {$r$};
      \fill[object] (2.1,2.0) circle (4.5pt);
      \node[diagramlabel, anchor=west, inner sep=8pt, font=\footnotesize] at (2.1,2.0) {stopper};
      \draw[hscmuted, line width=3pt] (0,0.95) -- (0,2.0);
      \node[diagramlabel, anchor=east, font=\footnotesize] at (-0.15,1.15) {tube};
      \draw[hscink, line width=0.7pt] (0,0.95) -- (0,0.45);
      \draw[hscink, line width=1.6pt] (-0.12,0.82) -- (0.12,0.82);
      \node[diagramlabel, anchor=west, font=\footnotesize] at (0.15,0.82) {marker};
      \draw[object] (-0.22,0.45) rectangle (0.22,-0.05);
      \node[diagramlabel, anchor=west, font=\footnotesize] at (0.3,0.2) {hanging mass};
      \draw[hscink, line width=0.6pt] (-2.4,-1.0) -- (2.6,-1.0);
      \draw[object] (-1.9,-0.95) rectangle (-1.1,-0.8);
      \draw[displacement] (-1.5,-0.75) -- (-1.5,-0.1);
      \node[diagramlabel, anchor=north, font=\footnotesize] at (-0.6,-1.05) {phone (films upwards)};
      \node[diagramlabel] at (0.1,-2.05) {\textbf{(a)} Part A};
    \end{tikzpicture}\hspace{12mm}%
    \begin{tikzpicture}[x=0.75cm, y=0.75cm]
      % ---- (b) Part B: pendulum and light gate
      \fill[ground] (-0.8,3.2) rectangle (0.8,3.42);
      \draw[hscink, line width=0.7pt] (-0.8,3.2) -- (0.8,3.2);
      \coordinate (O) at (0,3.2);
      \coordinate (L) at (0,0.4);
      \coordinate (S) at ($(O)+(-40:2.8)$);
      \draw[guide] (O) -- (L);
      \draw[hscink, line width=0.7pt] (O) -- (S);
      \draw[displacement, dashed] (S) arc (-40:-88:2.8);
      \draw[hscink!50, line width=0.8pt, dashed] (L) circle (4.5pt);
      \fill[object] (S) circle (4.5pt);
      \node[diagramlabel, anchor=south west, inner sep=4pt, font=\footnotesize] at (S) {released from rest};
      \draw[hscmuted, line width=2.6pt] (-0.42,0.8) -- (-0.42,0.02) -- (0.42,0.02) -- (0.42,0.8);
      \node[diagramlabel, anchor=south east, inner sep=2pt, font=\footnotesize] at (-0.45,0.8) {light gate};
      \draw[hscink, line width=0.5pt] (-0.42,0.3) -- (-1.25,0.3);
      \draw[object] (-3.35,0.02) rectangle (-1.25,0.63);
      \node[diagramlabel, font=\scriptsize] at (-2.3,0.325) {data logger};
      \draw[guide] ($(S)+(0.25,0)$) -- ($(S)+(1.0,0)$);
      \draw[guide] (0.55,0.4) -- ($(S|-L)+(1.0,0)$);
      \draw[hscink, <->, line width=0.5pt] ($(S)+(0.85,0)$) -- node[right, diagramlabel] {$h$} ($(S|-L)+(0.85,0)$);
      \node[diagramlabel] at (0,-1.0) {\textbf{(b)} Part B};
    \end{tikzpicture}
    \caption{(a) Part A: the stopper moves in a horizontal circle; the hanging mass keeps the
      tension the same. (b) Part B: the ball falls through height $h$.}\label{ch3:fig:pendulum}
  \end{figure}}
\method{\textbf{Part A (horizontal circle).}
\begin{enumerate}
  \item Thread the string through the tube. Tie the stopper to the top end and the hanging mass
        to the bottom end. The hanging mass keeps the tension the same.
  \item Hold the tube upright. Whirl the stopper in a horizontal circle of radius
        \qty{0.50}{m} at a steady rate. Keep the paper-clip marker just below the tube, so $r$
        stays the same.
  \item Film the circle from directly below. First film a metre rule at the height of the circle,
        for scale.
  \item Play the video frame by frame. At 8 points round the circle, measure the distance the
        stopper moves in 6 frames. 6 frames take $6 \div 240 = \qty{0.025}{s}$, so
        $v = \text{distance} \div \qty{0.025}{s}$. Find $K = \tfrac{1}{2}mv^2$ at each point. Do
        three runs.
\end{enumerate}
\textbf{Part B (vertical swing).}
\begin{enumerate}[resume]
  \item Release the ball from rest with its centre $h = \qty{0.050}{m}$ above its lowest point.
        The light gate is at the lowest point. Speed at the bottom $=$ diameter $\div$ time the
        beam is blocked. Change mm and ms to m and s first:
        $v = \qty{0.0250}{m} \div \qty{0.0256}{s} = \qty{0.977}{m.s^{-1}}$.
  \item Repeat three times, then for $h = 0.100$, 0.150 and \qty{0.200}{m}.
\end{enumerate}}
\results{\textbf{Part A} (sample results, one run; $m = \qty{0.025}{kg}$; distance moved in
  6 frames):
  \begin{center}\small
  \begin{tabular}{@{}lcccccccc@{}}
    \toprule
    Point & 1 & 2 & 3 & 4 & 5 & 6 & 7 & 8\\
    \midrule
    Distance (\si{cm}) & 11.0 & 11.1 & 10.9 & 11.0 & 11.1 & 11.0 & 10.9 & 11.0\\
    $v$ (\si{m.s^{-1}}) & 4.40 & 4.44 & 4.36 & 4.40 & 4.44 & 4.40 & 4.36 & 4.40\\
    $K$ (\si{J}) & 0.242 & 0.246 & 0.238 & 0.242 & 0.246 & 0.242 & 0.238 & 0.242\\
    \bottomrule
  \end{tabular}
  \end{center}
  The other two runs gave the same pattern. The height of the circle did not change.

  \textbf{Part B} (sample results, means of three trials; $m = \qty{0.0640}{kg}$):
  \begin{center}\small
  \begin{tabular}{@{}cccccc@{}}
    \toprule
    $h$ (\si{m}) & time (\si{ms}, milliseconds) & $v$ (\si{m.s^{-1}}) & $K = \tfrac{1}{2}mv^2$ (\si{J})
      & $mgh$ (\si{J}) & $K \div mgh$\\
    \midrule
    0.050 & 25.6 & 0.977 & 0.0305 & 0.0314 & 0.97\\
    0.100 & 18.1 & 1.381 & 0.0610 & 0.0627 & 0.97\\
    0.150 & 14.8 & 1.689 & 0.0913 & 0.0941 & 0.97\\
    0.200 & 12.8 & 1.953 & 0.122 & 0.125 & 0.97\\
    \bottomrule
  \end{tabular}
  \end{center}}
\conclusion{\textbf{Part A:} $K$ stayed at \qty{0.24}{J} at all 8 points. The small changes (up to
  2\,\%) go up and down at random: they come from reading the video to $\pm\qty{0.1}{cm}$. The height
  did not change, so $U$ was constant. So $E = K + U$ was constant, and the net work on the stopper
  was zero. The tension is at \ang{90} to the velocity, so it does no work. (Your hand does a
  little work to replace what air resistance takes.)

  \textbf{Part B:} $K$ at the bottom was 97\,\% of the $U$ lost, the same at every height. Air
  resistance on this small, heavy ball takes less than 0.2\,\% in one swing. So the missing
  3\,\% is a \term{systematic error} (an error that makes every reading wrong by the same
  fraction). For example, the beam is not a thin line. If the gate counts it as blocked while any
  part of the ball covers it, the ball moves about \qty{25.4}{mm} (not \qty{25.0}{mm}) while the
  beam is blocked. Then every speed is 1.6\,\% too low and every $K$ is 3\,\% too low. Within this
  error, gravity did all the work: the $K$ gained equals the $U$ lost.}
\significance{Energy conservation explains pendulums, rollercoasters and satellites. The work done
  by other forces (motors, friction) explains when the total energy changes.}
\validity{\begin{itemize}[beginpenalty=10000]
  \item \textbf{Validity:} A: keep the circle horizontal and $r$ constant (watch the marker), so
        $U$ stays constant. B: release from rest. Measure $h$ to the ball's centre.
  \item \textbf{Reliability:} A: three runs, 8 points in each. B: repeat each height three times
        and check that the times agree.
  \item \textbf{Accuracy:} use a light gate (not a stopwatch). Measure the diameter with vernier
        calipers. Line up the beam with the centre of the ball (check with the ball hanging at
        rest). Better still, push the ball slowly through the beam along a ruler. Note where the
        gate turns on and off, and use this blocked length instead of the diameter.
\end{itemize}}
\end{experiment}

% Added by the syllabus coverage audit (M5.IQ2.04 and M5.IQ2.05 say "investigate", but this
% section gave no model-answer points for "describe an investigation" questions). No numbers.
% Final audit: the last paragraph now covers both parts (horizontal circle: K constant; swing:
% gravity does work, K gained = U lost), and summary line 2 says "uniform circular motion".
\begin{examtip}
``Describe an investigation to show how torque depends on the distance from the pivot''
(3--4 marks). Make these points:
\begin{enumerate}
  \item Balance a metre rule on a pivot at its centre of mass, so its own weight has no torque.
  \item Hang a fixed load at a fixed distance on one side. Its torque stays the same.
  \item On the other side, pull down at \ang{90} to the rule with a spring balance. Record the
        force $F$ that keeps the rule level. Do this at several distances $r$, three times each.
  \item Plot $F$ against $1/r$. A straight line through the origin shows that $Fr$ is
        constant, as $\tau = r_\perp F$ predicts.
\end{enumerate}
For the work--energy investigation, link each result to the direction of the force. In the
horizontal circle, the tension is at \ang{90} to the velocity, so it does no work, and $K$ stays
constant. In the swing, the tension still does no work, but gravity does: the $K$ gained equals
the $U$ lost.
\end{examtip}

%% ==================================================================
% Chapter summary (added by the whole-book consistency edit: "How to use this guide" says
% that every chapter ends with a short summary). Every statement is taught above.
\begin{summarybox}[unbreakable]{\Cref{ch3:chap} in five lines}
\begin{itemize}
  \item Use radians: one turn is $2\pi$ rad. $\omega = \frac{\Delta\theta}{t} = \frac{2\pi}{T}$ and
        $v = \omega r$. Every point on a spinning object has the same $\omega$.
  \item In uniform circular motion the net (centripetal) force is at \ang{90} to the motion, so
        it does no work ($W = Fs\cos\ang{90} = 0$). So $K$ and the speed stay constant.
  \item Only forces other than gravity change the total energy:
        $W_\mathrm{other} = \Delta E = \Delta K + \Delta U$.
  \item Vertical circle: use energy for the speed (top to bottom is a drop of $2r$). Then use
        the forces at the top ($F_\mathrm{T} + mg = \frac{mv^2}{r}$) or the bottom
        ($F_\mathrm{T} - mg = \frac{mv^2}{r}$).
  \item Torque: $\tau = r_\perp F = rF\sin\theta$, with $r$ measured from the pivot. When the
        torques balance, the rotation does not change.
\end{itemize}
\end{summarybox}
